\documentclass[11pt]{article}
\usepackage[margin=25mm]{geometry}
\usepackage{fontspec}
\setmainfont{Latin Modern Roman}
\usepackage{amsmath,amssymb,mathtools}
\usepackage{unicode-math}
\setmathfont{Latin Modern Math}
\usepackage{xurl,hyperref}
\hypersetup{colorlinks=true,urlcolor=blue}
\setcounter{secnumdepth}{0}
\setlength{\parindent}{0pt}
\setlength{\parskip}{5pt}
\setlength{\emergencystretch}{3em}
\newcommand{\sref}[2]{\hyperlink{source:#1:#2}{[S#2]}}
\newcommand{\cataloguescope}[1]{\paragraph{Recorded target.}#1}
\newcommand{\R}{\mathbb R}
\newcommand{\C}{\mathbb C}
\newcommand{\N}{\mathbb N}
\newcommand{\Z}{\mathbb Z}
\newcommand{\Q}{\mathbb Q}
\newcommand{\F}{\mathbb F}
\newcommand{\eps}{\varepsilon}
\begin{document}
\section{Classification of Spatial and Nonspatial Markov Semigroups}\label{statement}

\subsection{Definitions and mathematical statement}

\paragraph{Shared notation.}$\mathbb N=\{1,2,\ldots\}$; the standard number fields have their usual meanings. All additional conventions are stated locally.


\paragraph{Definitions and conventions.} A Markov semigroup $T_t$ is a continuous semigroup of unital completely positive maps on the specified operator algebra $B$. A dilation consists of an embedding $i:B\to A$, a conditional expectation or compression $p:A\to B$, and endomorphisms $\vartheta_t$ with $T_t=p\vartheta_t i$. The GNS construction for $p$ gives a Hilbert $B$-module $E$ and a represented algebra in $\mathcal B^a(E)$. An automorphic dilation has invertible $*$-automorphisms; an inner one is implemented by unitaries. A spatial Markov semigroup has the appropriate dominating elementary CP-semigroup/unit structure of the source and an associated product system with a unit. A Fock or type-I product system is generated by its units; a spatial non-Fock system can have type II. Here nontrivial nonspatial excludes merely choosing a nonspatial $E_0$-semigroup or taking its tensor product with a spatial Markov semigroup. A form generator defines the semigroup infinitesimally through quadratic or sesquilinear forms, potentially with unbounded operator data.

\paragraph{Source scope and qualifications.} The source explicitly excludes trivial nonspatial endomorphism and tensor-product examples. Its complete paragraph also asks which semigroups on a subalgebra admit spatial extensions and how such an extension changes product-system type; these associated questions are retained as source context.

\paragraph{Further source questions.} For a von Neumann algebra $B\subseteq B(H)$, determine which nonspatial Markov semigroups admit spatial extensions to $B(H)$, and how an embedding into the product system of an extension changes the type classification.

\paragraph{Catalogue mathematical target.}
\paragraph{Catalogue wording retained for traceability.}
For Markov semigroups, determine whether nontrivial nonspatial examples exist on $B(H)$; whether spatial Markov semigroups on $B(H)$ can have non-Fock product systems; whether all spatial examples are known or further ones exist; whether such further examples could yield non-type-I Markov semigroups; and whether the theory of form generators can be adapted from $B(H)$ to general von Neumann algebras.

\paragraph{Saved literature context.} The supplied catalogue's assessment dated 2026-08-21 records: Later work proves Hudson--Parthasarathy dilations for spatial Markov semigroups on B(H), but the listed classification of nonspatial examples, non-Fock product systems, and general von Neumann algebra extensions remains incomplete. The assessment is retained as dated provenance; it is not a new comprehensive literature or open-status certification.
\subsection{Short English statement}
Classify genuinely nonspatial Markov examples, spatial examples with non-Fock product systems, and the reach of form-generator methods.
\subsection{Research attempt}
\paragraph{Model, date and investigated scope.}
OpenAI GPT-6 Astra (configured model \texttt{gpt-6-astra}), ultra reasoning; generated 11 October 2026. Proves CP domination by an elementary semigroup for an explicitly bounded generator ansatz, explaining why it cannot generate the requested nonspatial examples.

\paragraph{Sources and literature check.}
The complete supplied definition, including source qualifications and the dated literature assessment, was reviewed. The calculation below is an elementary restricted attempt motivated by [S2]; it does not claim novelty or a fresh literature-status determination.

\paragraph{Context and investigated scope.} The complete supplied definition, including source qualifications and the dated literature assessment, was reviewed. The calculation below is an elementary restricted attempt motivated by [S2]; it does not claim novelty or a fresh literature-status determination.

\paragraph{Why a bounded-generator search stays spatial.} Consider on $B(\mathcal H)$ a bounded generator of the explicitly assumed form
\[
 L(X)=K^*X+XK+\Phi(X),
\]
where $K$ is bounded, $\Phi$ is a bounded completely positive map, and $L(I)=0$. Put $S_t(X)=e^{tK^*}Xe^{tK}$. Variation of constants gives the norm-convergent Dyson expansion for $T_t=e^{tL}$ as $S_t$ plus integrals of compositions $S_{t-s_r}\Phi S_{s_r-s_{r-1}}\cdots\Phi S_{s_1}$ over ordered times. Each composition and its integral is completely positive. The remainder norm is bounded by an exponential-series tail because all generators are bounded.

\paragraph{An elementary CP lower bound.} It follows that $T_t-S_t$ is completely positive for every $t\geq0$. The operators $e^{tK}$ form a semigroup, so $S_t$ is the elementary CP semigroup required by the relevant domination notion of spatiality. Hence this entire bounded ansatz supplies spatial examples, not the source's sought genuinely nonspatial Markov examples. No theorem that every generator has this form is needed for this limited conclusion.

\paragraph{What the calculation leaves open.} To escape this argument one must leave the assumed bounded decomposition, for instance through domain-sensitive unbounded form generators. That does not by itself produce a nonspatial example. The associated product-system type, possible spatial non-Fock examples, extension from a subalgebra to $B(\mathcal H)$ and adaptation of form methods to general von Neumann algebras are all unaddressed. Trivial tensor products with nonspatial endomorphism semigroups, expressly excluded in the source, are not proposed as answers.

\paragraph{Verification and adversarial audit.}
Checked ordered Dyson terms, norm convergence and CP domination; did not invoke an unchecked universal generator representation theorem. Model self-audit only; no independent human proof checking. Prompts A and B were applied to this canonical record and its qualified source scope. Checked ordered Dyson terms, norm convergence and CP domination; did not invoke an unchecked universal generator representation theorem. Retained unresolved status and zero original-target progress; Unbounded-domain constructions and product-system/extension analysis needed for the bundled classification questions are missing.

Coordinator checked the completely positive Dyson terms and elementary semigroup domination under the explicitly bounded generator decomposition. This additional check was model-assisted, not independent human verification.

\paragraph{Tools, process and contributions.}
Python 3 was used for file preparation and structural checks; no numerical experiment or theorem prover was used. Contribution prompt A was followed by prompt B. GPT-6 Astra with ultra reasoning drafted the research and performed its stated self-audit; any additional coordinator contribution is identified in the verification and audit. A GPT-6-based Codex coordinator assembled the manuscript and checked repository integration; its exact variant and reasoning configuration were not exposed. No human mathematical contribution, independent human proof check or formal proof certification is supplied. Complete derivations and the consulted source references are included here; private reasoning transcripts are not published.

\paragraph{Final status and first remaining gap.}
\textbf{UNRESOLVED}. Unbounded-domain constructions and product-system/extension analysis needed for the bundled classification questions are missing.

\section{Completion Estimate}
\noindent\textbf{COMPLETION: 0\%}
The calculation is an auxiliary or restricted attempt and establishes no newly verified portion of the full source target. Unbounded-domain constructions and product-system/extension analysis needed for the bundled classification questions are missing.

\subsection{Sources}
\begin{itemize}
\item[\hypertarget{source:statement:1}{S1}] ULAM.AI and the UnsolvedMath Contributors, UnsolvedMath, catalogue entry. Catalogue curation under CC BY 4.0; definitions and English exposition added locally.
\url{https://huggingface.co/datasets/ulamai/UnsolvedMath}.
\item[\hypertarget{source:statement:2}{S2}] Mini-Workshop: Product Systems and Independence in Quantum Dynamics (2009), Oberwolfach Reports, DOI 10.4171/owr/2009/09, pp. 493--548; official report source and mathematical context.
\url{https://doi.org/10.4171/owr/2009/09}.
\item[\hypertarget{source:statement:3}{S3}] Official EMS Press PDF of the same Oberwolfach report; PDF page 52, printed page 544.
\url{https://ems.press/content/serial-article-files/46209}.
\item[\hypertarget{source:statement:4}{S4}] spatial Markov semigroups admit dilations; reference in the saved dated catalogue literature review.
\url{https://doi.org/10.3842/SIGMA.2022.071}.
\item[\hypertarget{source:statement:5}{S5}] spatial CP-semigroups and product systems; reference in the saved dated catalogue literature review.
\url{https://arxiv.org/abs/0804.2169}.
\end{itemize}
\label{end:statement}
\end{document}
